\section{Related Work}
\label{sec:related}

\subsection{Fretting Corrosion and Connector Reliability}
The physics of electrical contacts is mature~\cite{holm,braunovic}. Under
small-amplitude relative motion, the oxide-free metallic spots of a
tin-plated contact are progressively replaced by an insulating debris layer;
contact resistance stays near its initial milliohm level for an incubation
period and then rises steeply, often by several orders of magnitude, within
a comparatively small number of additional cycles~\cite{antler1985,
bryant1994,park2008}. Vibration-driven fretting in automotive connectors,
and the dependence of the degradation rate on vibration amplitude and
spectrum, have been studied experimentally~\cite{flowers2004}. Industry
qualification tests (e.g., USCAR-2~\cite{uscar2}) define an electrical
discontinuity as a resistance excursion above \SI{7}{\ohm} lasting at least
\SI{1}{\micro\second} during vibration. These works characterize the
\emph{connector}; they do not address how a degraded connector manifests in
the \emph{data} of the system it serves, which is the question here.

\subsection{Timing and Physical-Layer Analysis of In-Vehicle Networks}
A large literature exploits the timing and electrical characteristics of
Controller Area Network (CAN) traffic~\cite{bosch_can,iso11898}, mostly for
intrusion detection. Inter-arrival-time analysis of periodic
messages~\cite{song2016}, clock-skew fingerprinting of ECUs~\cite{cho_cids}
and voltage fingerprinting~\cite{murvay2014,cho_viden,voltageids,scission,
ecuprint} identify senders or detect injected frames; deep learning
detectors operate on frame sequences~\cite{kang2016,taylor2016,canet}.
Public datasets such as HCRL Car-Hacking~\cite{hcrl_gids},
ROAD~\cite{road}, can-train-and-test~\cite{cantrain} and
CAN-MIRGU~\cite{mirgu} serve this community. We borrow the timing view, but
our target is the opposite of an attack: frames that \emph{should} arrive
but do not, and only at one observation point. To our knowledge no prior
work uses in-vehicle network telemetry to assess the health of the
diagnostic connector itself.

\subsection{Diagnostics, Anomaly Detection and Prognostics}
Model-based fault detection by analytical redundancy~\cite{chow1984} and
data-driven anomaly detection for automotive systems~\cite{theissler2017}
are established, as are general prognostics and health management (PHM)
methods for estimating RUL~\cite{si2011,lei2018} and their evaluation
metrics~\cite{saxena2010}. Unsupervised detectors such as Isolation
Forest~\cite{iforest}, one-class SVM~\cite{ocsvm} and LSTM
encoder--decoders~\cite{malhotra2016} are standard baselines and are used as
such here. What distinguishes the present problem is that (i) the degrading
component has no sensor of its own, (ii) its symptoms are shared with
several benign or unrelated causes, and (iii) the stimulus that drives the
degradation is itself measurable in the same data stream. ECTA is designed
around these three properties.
